% ══════════════════════════════════════════════════════════════════
% Section 6 — Benchmark Results
% ══════════════════════════════════════════════════════════════════
\section{Benchmark results}
\label{sec:results}

All benchmarks were run using the scripts in \texttt{paper/run\_benchmarks.py}
with 10\,000 training samples, 1000 epochs, and early stopping with
patience~200.

\subsection{Accuracy comparison}
\label{sec:results:accuracy}

\Cref{tab:accuracy} reports the coefficient of determination ($R^2$) for
energy and stress (invariant gradients) across all material--architecture
combinations.

% ── Main accuracy table ───────────────────────────────────────────
\begin{table}[htbp]
\centering
\caption{Surrogate accuracy ($R^2$) for energy and stress predictions.  Best
  result per material in \textbf{bold}.}
\label{tab:accuracy}
% TODO: Replace with actual results from results/benchmarks.tex or
%   auto-generated LaTeX table.
\begin{tabular}{@{}ll S[table-format=1.4] S[table-format=1.4] r r@{}}
\toprule
\textbf{Material} & \textbf{Model}
  & {$R^2$ (energy)} & {$R^2$ (stress)}
  & {\textbf{Params}} & {\textbf{Time (s)}} \\
\midrule
\multirow{3}{*}{Neo-Hookean}
  & MLP-64x64       & {[X.XXXX]} & {[X.XXXX]} & {[N]} & {[T]} \\
  & ICNN-64x64      & {[X.XXXX]} & {[X.XXXX]} & {[N]} & {[T]} \\
  & Polyconvex-64x64& {[X.XXXX]} & {[X.XXXX]} & {[N]} & {[T]} \\
\midrule
\multirow{3}{*}{Mooney--Rivlin}
  & MLP-64x64       & {[X.XXXX]} & {[X.XXXX]} & {[N]} & {[T]} \\
  & ICNN-64x64      & {[X.XXXX]} & {[X.XXXX]} & {[N]} & {[T]} \\
  & Polyconvex-64x64& {[X.XXXX]} & {[X.XXXX]} & {[N]} & {[T]} \\
\midrule
\multirow{3}{*}{Yeoh}
  & MLP-64x64       & {[X.XXXX]} & {[X.XXXX]} & {[N]} & {[T]} \\
  & ICNN-64x64      & {[X.XXXX]} & {[X.XXXX]} & {[N]} & {[T]} \\
  & Polyconvex-64x64& {[X.XXXX]} & {[X.XXXX]} & {[N]} & {[T]} \\
\midrule
\multirow{2}{*}{Holzapfel--Ogden}
  & MLP-64x64       & {[X.XXXX]} & {[X.XXXX]} & {[N]} & {[T]} \\
  & ICNN-64x64      & {[X.XXXX]} & {[X.XXXX]} & {[N]} & {[T]} \\
\midrule
\multirow{2}{*}{GOH}
  & MLP-64x64       & {[X.XXXX]} & {[X.XXXX]} & {[N]} & {[T]} \\
  & ICNN-64x64      & {[X.XXXX]} & {[X.XXXX]} & {[N]} & {[T]} \\
\bottomrule
\end{tabular}
\end{table}

% ── Accuracy bar charts ──────────────────────────────────────────
\begin{figure}[htbp]
\centering
\begin{subfigure}[t]{0.48\textwidth}
  % TODO: Include figures/accuracy_comparison_energy.pdf
  \fbox{\parbox{\textwidth}{\centering\vspace{4cm}
    \textbf{[PLACEHOLDER]} \\
    \texttt{figures/accuracy\_comparison\_energy.pdf}
  \vspace{1cm}}}
  \caption{$R^2$ for energy prediction.}
  \label{fig:accuracy_energy}
\end{subfigure}
\hfill
\begin{subfigure}[t]{0.48\textwidth}
  % TODO: Include figures/accuracy_comparison_stress.pdf
  \fbox{\parbox{\textwidth}{\centering\vspace{4cm}
    \textbf{[PLACEHOLDER]} \\
    \texttt{figures/accuracy\_comparison\_stress.pdf}
  \vspace{1cm}}}
  \caption{$R^2$ for stress (invariant gradient) prediction.}
  \label{fig:accuracy_stress}
\end{subfigure}
\caption{Accuracy comparison across materials and architectures.}
\label{fig:accuracy_comparison}
\end{figure}

\subsection{Convergence behavior}
\label{sec:results:convergence}

\Cref{fig:convergence} shows the validation loss during training for all
material--model combinations.

% ── Convergence curves ────────────────────────────────────────────
\begin{figure}[htbp]
\centering
% TODO: Include figures/convergence_curves.pdf
\fbox{\parbox{0.9\textwidth}{\centering\vspace{5cm}
  \textbf{[PLACEHOLDER]} \\
  \texttt{figures/convergence\_curves.pdf} \\[0.5em]
  Grid of validation loss curves (log scale) for each material, with one
  curve per architecture.  Shows convergence speed differences between MLP,
  ICNN, and Polyconvex.
\vspace{1cm}}}
\caption{Validation loss convergence for all material--architecture
  combinations.  The ICNN and polyconvex models typically require more
  epochs to converge due to constrained weight updates.}
\label{fig:convergence}
\end{figure}

\subsection{Stress prediction quality}
\label{sec:results:stress}

To assess the quality of stress predictions derived from the energy surrogate
via automatic differentiation, \Cref{fig:stress_scatter} shows scatter plots
of predicted versus analytical invariant gradients $\partial\sef/\partial
I_\alpha$ on the validation set for the Neo-Hookean material.

% ── Stress scatter ────────────────────────────────────────────────
\begin{figure}[htbp]
\centering
% TODO: Include figures/stress_scatter.pdf
\fbox{\parbox{0.9\textwidth}{\centering\vspace{4.5cm}
  \textbf{[PLACEHOLDER]} \\
  \texttt{figures/stress\_scatter.pdf} \\[0.5em]
  Predicted vs.\ analytical $\partial\sef/\partial I_\alpha$ for MLP and
  ICNN on Neo-Hookean.  Points color-coded by invariant component
  ($\partial W/\partial I_1$, $\partial W/\partial I_2$, $\partial W/\partial J$).
  Ideal predictions lie on the diagonal.
\vspace{1cm}}}
\caption{Predicted versus analytical invariant gradients (Neo-Hookean,
  validation set).  Both MLP and ICNN achieve tight clustering around the
  identity line.}
\label{fig:stress_scatter}
\end{figure}

\subsection{Error distribution}
\label{sec:results:error}

\Cref{fig:error_dist} shows the distribution of pointwise relative errors
for the invariant gradients.

% ── Error distribution ────────────────────────────────────────────
\begin{figure}[htbp]
\centering
% TODO: Include figures/error_distribution.pdf
\fbox{\parbox{0.7\textwidth}{\centering\vspace{4cm}
  \textbf{[PLACEHOLDER]} \\
  \texttt{figures/error\_distribution.pdf} \\[0.5em]
  Histogram of relative errors in $\partial\sef/\partial I_\alpha$ for MLP
  and ICNN.  Legend shows median relative error per architecture.
\vspace{1cm}}}
\caption{Distribution of relative errors in invariant gradient predictions
  (Neo-Hookean, validation set).}
\label{fig:error_dist}
\end{figure}

\subsection{Scaling studies}
\label{sec:results:scaling}

We investigate how surrogate accuracy depends on the training set size and
network capacity.

\subsubsection{Training set size}

\Cref{fig:scaling_samples} shows the stress $R^2$ as a function of the number
of training samples for a fixed MLP-$64\times64$ architecture on Neo-Hookean.

% ── Scaling: samples ──────────────────────────────────────────────
\begin{figure}[htbp]
\centering
% TODO: Include figures/scaling_samples.pdf
\fbox{\parbox{0.6\textwidth}{\centering\vspace{4cm}
  \textbf{[PLACEHOLDER]} \\
  \texttt{figures/scaling\_samples.pdf} \\[0.5em]
  Semi-log plot of $R^2$(stress) vs.\ number of training samples.
  Shows diminishing returns beyond a threshold.
\vspace{1cm}}}
\caption{Surrogate accuracy versus training set size (MLP-$64\times64$,
  Neo-Hookean).}
\label{fig:scaling_samples}
\end{figure}

\subsubsection{Network width}

\Cref{fig:scaling_width} shows the accuracy and parameter count as the hidden
layer width varies from 16 to 256.

% ── Scaling: width ────────────────────────────────────────────────
\begin{figure}[htbp]
\centering
% TODO: Include figures/scaling_width.pdf
\fbox{\parbox{0.6\textwidth}{\centering\vspace{4cm}
  \textbf{[PLACEHOLDER]} \\
  \texttt{figures/scaling\_width.pdf} \\[0.5em]
  Dual-axis plot: $R^2$(stress) on left axis, parameter count (bar) on
  right axis, versus hidden layer width.
\vspace{1cm}}}
\caption{Surrogate accuracy versus network width (Neo-Hookean).  The bar
  chart shows the corresponding number of trainable parameters.}
\label{fig:scaling_width}
\end{figure}

\subsection{Inference speed}
\label{sec:results:timing}

\Cref{fig:timing} compares the wall-clock time for evaluating the stress
tensor at
% TODO: Insert actual n_eval from timing.json
\textbf{[N]} integration points using the analytical material model versus
the trained MLP surrogate (CPU, single-threaded).

% ── Timing bar chart ─────────────────────────────────────────────
\begin{figure}[htbp]
\centering
% TODO: Include figures/timing_bar.pdf
\fbox{\parbox{0.5\textwidth}{\centering\vspace{4cm}
  \textbf{[PLACEHOLDER]} \\
  \texttt{figures/timing\_bar.pdf} \\[0.5em]
  Bar chart comparing analytical vs.\ NN inference time (ms).
  Annotated with speedup factor.
\vspace{1cm}}}
\caption{Inference time comparison between the analytical Neo-Hookean model
  and the MLP-$64\times64$ surrogate.  Error bars show standard deviation
  over
  % TODO: Insert n_repeats.
  \textbf{[N]} repetitions.}
\label{fig:timing}
\end{figure}
